% ============================================================
% Section 94: 热发射电子的平均能量
% ============================================================
\section{固体理论：热发射电子的平均能量}
\label{sec:thermionic-energy}

\begin{modelbox}[题目：热发射电子垂直与平行表面方向的平均能量]
证明热发射电子在离开金属表面后，其垂直于金属表面的平均能量为 $k_BT$，平行于金属表面的平均能量也是 $k_BT$。
\end{modelbox}

\begin{tipbox}[考点快速分析]
\begin{itemize}
    \item \textbf{麦克斯韦-玻尔兹曼近似}：高温下 $k_BT\ll W$，费米分布尾端可用经典分布
    \item \textbf{逸出条件}：$E_\perp=p_z^2/2m\ge\chi=W+E_F$
    \item \textbf{通量积分}：对动量空间积分，计算总粒子通量和能量通量
    \item \textbf{扣除势垒}：离开表面后的能量需减去 $\chi$
\end{itemize}
\end{tipbox}

\begin{derivationbox}[标准解题过程]
\textbf{基本模型}

金属表面位于 $z=0$，内部 $z<0$。电子被深度为 $\chi$ 的势阱束缚。功函数
\begin{equation}
W=\chi-E_F.
\end{equation}

电子逸出条件：垂直动能 $E_\perp=p_z^2/2m\ge\chi$，即 $p_z\ge p_{z0}=\sqrt{2m\chi}$。

高温近似（$k_BT\ll W$）下，电子服从麦克斯韦-玻尔兹曼分布
\begin{equation}
f(\bm{p})\approx\exp\!\Bigl[-\frac{E(\bm{p})-E_F}{k_BT}\Bigr].
\end{equation}

单位时间、单位面积、动量在 $d^3p$ 范围内且能逸出的电子数
\begin{equation}
d\nu=v_z\,dn=\frac{2}{mh^3}p_z\exp\!\Bigl[-\frac{E-E_F}{k_BT}\Bigr]d^3p.
\end{equation}

\textbf{Step 1: 垂直方向的平均能量}

总逸出电子通量（理查德森-杜什曼公式）
\begin{equation}
N=\int d\nu=\frac{4\pi m(k_BT)^2}{h^3}e^{-W/k_BT}.
\end{equation}

垂直能量通量
\begin{equation}
\Phi_\perp=\int\frac{p_z^2}{2m}\,d\nu
=\frac{e^{E_F/k_BT}}{m^2h^3}\int_{p_{z0}}^\infty p_z^3e^{-p_z^2/2mk_BT}dp_z
\iint_{-\infty}^\infty e^{-(p_x^2+p_y^2)/2mk_BT}dp_xdp_y.
\end{equation}

计算积分：
\begin{itemize}
    \item 二维高斯积分：$\iint e^{-(p_x^2+p_y^2)/2mk_BT}dp_xdp_y=2\pi mk_BT$
    \item $p_z$ 积分（令 $u=p_z^2$）：$\int_{p_{z0}^2}^\infty ue^{-u/2mk_BT}du=(2mk_BT)^2e^{-\chi/k_BT}(\chi/k_BT+1)$
\end{itemize}

代入得
\begin{equation}
\Phi_\perp=\frac{4\pi m(k_BT)^3}{h^3}e^{-W/k_BT}(k_BT+\chi).
\end{equation}

逸出电子在金属\textbf{内部}的平均垂直能量
\begin{equation}
\bar E_\perp=\frac{\Phi_\perp}{N}=k_BT+\chi.
\end{equation}

离开表面后，需扣除克服势垒的能量 $\chi$，故\textbf{外部}平均垂直动能
\begin{equation}
\boxed{\bar E_\perp^{\text{out}}=\bar E_\perp-\chi=k_BT.}
\end{equation}

\textbf{Step 2: 平行方向的平均能量}

平行能量通量
\begin{equation}
\Phi_\parallel=\int\frac{p_x^2+p_y^2}{2m}\,d\nu
=\frac{e^{E_F/k_BT}}{m^2h^3}\int_{p_{z0}}^\infty p_z e^{-p_z^2/2mk_BT}dp_z
\iint(p_x^2+p_y^2)e^{-(p_x^2+p_y^2)/2mk_BT}dp_xdp_y.
\end{equation}

计算积分：
\begin{itemize}
    \item $p_z$ 积分：$\int_{p_{z0}}^\infty p_z e^{-p_z^2/2mk_BT}dp_z=mk_BT e^{-\chi/k_BT}$
    \item 二维积分：$\iint(p_x^2+p_y^2)e^{-(p_x^2+p_y^2)/2mk_BT}dp_xdp_y=4\pi m^2(k_BT)^2$
\end{itemize}

代入得
\begin{equation}
\Phi_\parallel=\frac{2\pi m(k_BT)^3}{h^3}e^{-W/k_BT}.
\end{equation}

平行方向平均能量
\begin{equation}
\boxed{\bar E_\parallel=\frac{\Phi_\parallel}{N}=\frac{2\pi m(k_BT)^3}{4\pi m(k_BT)^2}=k_BT.}
\end{equation}
\end{derivationbox}

\begin{formulabox}[答案汇总]
\begin{center}
\begin{tabular}{ll}
\toprule
\textbf{方向} & \textbf{离开表面后的平均动能} \\
\midrule
垂直于表面 & $\bar E_\perp^{\text{out}}=k_BT$ \\
平行于表面 & $\bar E_\parallel=k_BT$ \\
\bottomrule
\end{tabular}
\end{center}
\end{formulabox}

\begin{insightbox}[物理图像]
\begin{itemize}
    \item 热发射电子的能量分布是金属内部费米-狄拉克分布的高能尾端，可近似为麦克斯韦-玻尔兹曼分布。每个自由度贡献 $\tfrac12k_BT$ 的动能，但垂直方向额外需要克服势垒 $\chi$。
    \item 离开表面后，垂直方向的平均动能恰好为 $k_BT$（扣除 $\chi$ 后），平行方向也为 $k_BT$。两者之和 $2k_BT$ 对应两个横向自由度的热运动能量。
    \item 这一结果是热电子发射理论和理查德森-杜什曼公式的基础，也是真空电子器件（电子显微镜、X射线管、真空二极管）中电子束初始能量分布的理论依据。
\end{itemize}
\end{insightbox}
